\begin{statementbox}
	\textbf{\textcolor{CStatement}{条件}}
	\begin{description}[style=nextline, leftmargin=2.8em, labelsep=0.5em, itemsep=0.35em]
		\item[\textbf{\textcolor{CStatement}{条件1（方向不平行）：}}] 直线$l_1$过点$A$，方向向量$\vec{u}$；直线$l_2$过点$B$，方向向量$\vec{v}$，且$\vec{u}\nparallel\vec{v}$（即$\vec{u}\times\vec{v}\neq\vec{0}$）。
	\end{description}

	\textbf{\textcolor{CStatement}{结论}}
	\begin{description}[style=nextline, leftmargin=2.8em, labelsep=0.5em, itemsep=0.45em]
		\item[\textbf{\textcolor{CStatement}{结论一（异面距离公式）：}}]
			\textbf{\textcolor{CStatement}{直观描述：}}
			两异面直线的距离等于向量$\overrightarrow{AB}$在公垂线方向投影的绝对值。
			\[
				d = \frac{\bigl| \overrightarrow{AB} \cdot (\vec{u} \times \vec{v}) \bigr|}{|\vec{u} \times \vec{v}|}.
			\]

		\item[\textbf{\textcolor{CStatement}{推广结论（法向量表示）：}}]
			\textbf{\textcolor{CStatement}{直观描述：}}
			记公垂线方向向量$\vec{n} = \vec{u}\times\vec{v}$，则距离为$\overrightarrow{AB}$在$\vec{n}$上投影长度。
			\[
				d = \frac{|\overrightarrow{AB} \cdot \vec{n}|}{|\vec{n}|}, \quad \vec{n} = \vec{u}\times\vec{v}.
			\]
	\end{description}

	\medskip
	\textbf{\textcolor{CStatement}{适用范围：}}
	公式用于计算异面直线距离，前提是两直线不平行（$\vec{u}\times\vec{v}\neq\vec{0}$）；若直线平行（$\vec{u}\times\vec{v}=\vec{0}$），公式分母为零，应改用点到直线距离公式。
\end{statementbox}
